% ============================================================
%  AutoPipeAD: Selección Automática de Pipelines de
%  Preprocesamiento para Detección de Anomalías en Imágenes
%  Artículo Técnico — CE-1114, I Semestre 2026
%  Instituto Tecnológico de Costa Rica
%  Formato IEEE Conference
% ============================================================
\documentclass[conference,spanish]{IEEEtran}

% ----- Codificación y lenguaje -----
\usepackage[utf8]{inputenc}
\usepackage[T1]{fontenc}
\usepackage[spanish,es-nodecimaldot]{babel}

% ----- Matemáticas -----
\usepackage{amsmath}
\usepackage{amssymb}

% ----- Tablas y figuras -----
\usepackage{graphicx}
\usepackage{booktabs}
\usepackage{array}
\usepackage{multirow}

% ----- Referencias y URLs -----
\usepackage{url}
\usepackage{cite}

% ----- Tipografía -----
\usepackage{microtype}

% ----- Código fuente -----
\usepackage{listings}
\lstset{
    basicstyle=\scriptsize\ttfamily,
    breaklines=true,
    breakatwhitespace=false,
    columns=flexible,
    keepspaces=true,
    showstringspaces=false,
    aboveskip=4pt,
    belowskip=2pt,
}

% ----- Diagramas TikZ -----
\usepackage{xcolor}
\usepackage{tikz}
\usetikzlibrary{shapes.geometric, arrows.meta, positioning, fit, backgrounds, calc}

% ----- Hiperenlaces (cargado al final para compatibilidad con IEEEtran) -----
\usepackage[hidelinks,breaklinks]{hyperref}

% ============================================================
%  COLORES
% ============================================================
\definecolor{tecblue}{RGB}{0,51,102}
\definecolor{lightblue}{RGB}{210,230,255}
\definecolor{lightgreen}{RGB}{210,240,220}
\definecolor{lightyellow}{RGB}{255,250,210}
\definecolor{lightgray}{RGB}{240,240,240}
\definecolor{lightred}{RGB}{255,220,220}

% ============================================================
%  ESTILOS TIKZ
% ============================================================
\tikzstyle{modulo}   = [rectangle, rounded corners=3pt, minimum width=2.2cm,
    minimum height=0.85cm, text centered, draw=tecblue, fill=lightblue,
    font=\scriptsize\bfseries, text width=2.0cm]
\tikzstyle{externo}  = [rectangle, minimum width=2.0cm,
    minimum height=0.85cm, text centered, draw=gray, fill=lightgray,
    font=\scriptsize, text width=1.85cm, dashed]
\tikzstyle{actor}    = [rectangle, minimum width=2.0cm,
    minimum height=0.85cm, text centered, draw=tecblue, fill=lightyellow,
    font=\scriptsize, text width=1.85cm]
\tikzstyle{estado}   = [rectangle, rounded corners=4pt, minimum width=2.6cm,
    minimum height=0.8cm, text centered, draw=tecblue, fill=lightgreen,
    font=\scriptsize, text width=2.4cm]
\tikzstyle{decision} = [diamond, aspect=2.2, minimum width=2.2cm,
    minimum height=0.7cm, text centered, draw=tecblue, fill=lightyellow,
    font=\scriptsize, text width=2.0cm]
\tikzstyle{startend} = [circle, minimum size=0.5cm, draw=tecblue, fill=tecblue]
\tikzstyle{iface}    = [rectangle, rounded corners=2pt, minimum width=2.2cm,
    minimum height=0.85cm, text centered, draw=tecblue!60, fill=lightred,
    font=\scriptsize\itshape, text width=2.0cm]
\tikzstyle{flecha}      = [-{Stealth[length=5pt]}, thick, color=tecblue]
\tikzstyle{flechadash}  = [-{Stealth[length=5pt]}, dashed, color=gray]
\tikzstyle{flechaopc}   = [-{Stealth[length=5pt]}, dotted, thick, color=tecblue!60]

% ============================================================
%  METADATOS PDF
% ============================================================
\hypersetup{
    pdftitle={AutoPipeAD: Selección Automática de Pipelines de Preprocesamiento
              para Detección de Anomalías en Imágenes},
    pdfauthor={Ricardo Antonio Borbón Mena}
}

% ============================================================
\begin{document}

\title{AutoPipeAD: Selección Automática de Pipelines de Preprocesamiento
       para Detección de Anomalías en Imágenes}

\author{%
\IEEEauthorblockN{Ricardo Antonio Borbón Mena}
\IEEEauthorblockA{Escuela de Ingeniería en Computadores\\
Instituto Tecnológico de Costa Rica\\
Cartago, Costa Rica\\
ricardobm@estudiantec.cr}
\and
\IEEEauthorblockN{Luis Alberto Chavarría Zamora \textit{(Supervisor)}}
\IEEEauthorblockA{Escuela de Ingeniería en Computadores\\
Instituto Tecnológico de Costa Rica\\
Cartago, Costa Rica\\
lachavarria@itcr.ac.cr}
}

\maketitle

% ============================================================
%  RESUMEN
% ============================================================
\begin{abstract}
La detección de anomalías en imágenes es un problema relevante en inspección
industrial, control de calidad y análisis médico. El desempeño de los
detectores de anomalías depende en gran medida del preprocesamiento aplicado
a las imágenes de entrada, y elegir manualmente esta cadena de
transformaciones toma tiempo y requiere conocimiento especializado en visión
computacional. En este trabajo se presenta AutoPipeAD, una librería Python de
código abierto que automatiza la búsqueda del pipeline de preprocesamiento
óptimo para un modelo base de detección de anomalías provisto por el usuario.
La librería recibe conjuntos de imágenes organizados en subdirectorios de
muestras normales y anómalas, y evalúa sistemáticamente combinaciones de
siete transformaciones predefinidas mediante
búsqueda exhaustiva por cuadrícula u optimización bayesiana con Optuna
(muestreo TPE y poda por mediana). La evaluación paralela se
implementó con joblib. Entre las métricas
configurables se incluyen AUC-ROC, F1-score, precisión, recall y tasa de
falsos positivos; el pipeline ganador se guarda en JSON o YAML para
garantizar reproducibilidad. Los experimentos con tres datasets sintéticos comprobaron el correcto funcionamiento de la librería y su compatibilidad con modelos base de scikit-learn y PyTorch.
\end{abstract}

\begin{IEEEkeywords}
detección de anomalías, preprocesamiento de imágenes, selección automática de
pipelines, búsqueda en cuadrícula, optimización bayesiana, Optuna, procesamiento
paralelo, IsolationForest, AUC-ROC, MVTec AD
\end{IEEEkeywords}

% ============================================================
%  I. INTRODUCCIÓN
% ============================================================
\section{Introducción}

La detección de anomalías en imágenes consiste en identificar
automáticamente instancias que se alejan de un patrón esperado, llamado
\emph{normal}~\cite{chandola2009anomaly}. Sus aplicaciones van desde la
inspección de defectos en superficies industriales hasta el análisis de
imágenes médicas y la vigilancia visual. A diferencia de la clasificación
supervisada, los detectores de anomalías suelen entrenarse asumiendo que las
muestras anómalas son escasas o desconocidas: el modelo aprende la
distribución de los casos normales y asigna una puntuación de anomalía a
cada instancia nueva~\cite{chandola2009anomaly}.

El desempeño de estos detectores depende en buena medida del
preprocesamiento aplicado a las imágenes de entrada. Operaciones como el
redimensionamiento, el suavizado gaussiano, la normalización estadística o
la ecualización adaptativa del histograma cambian la distribución de las
características de la imagen, y pueden mejorar o empeorar notablemente la
capacidad del modelo para distinguir lo normal de lo
anómalo~\cite{gonzalez2018dip}. Elegir esta cadena de transformaciones
—el \emph{pipeline de preprocesamiento}— se hace habitualmente a mano, por
prueba y error, lo que toma tiempo y exige experiencia en visión
computacional.

AutoPipeAD se desarrolló para automatizar este proceso. La librería recibe
un dataset de imágenes y un modelo base de detección de anomalías
compatible con la interfaz de scikit-learn o PyTorch, y busca
automáticamente la combinación de transformaciones que maximiza el
desempeño del modelo sobre un conjunto de validación.

El alcance del proyecto se limita a la selección del pipeline de
preprocesamiento: no incluye el entrenamiento ni la optimización de la
arquitectura del modelo base, ni clasificación multiclase, segmentación
semántica, detección de objetos, interfaz gráfica o despliegue como
servicio web.

% ============================================================
%  II. ANTECEDENTES
% ============================================================
\section{Antecedentes}

\subsection{Detección de Anomalías en Imágenes}

Chandola, Banerjee y Kumar~\cite{chandola2009anomaly} hicieron una revisión
amplia de los métodos de detección de anomalías, clasificándolos según el
nivel de supervisión disponible. En el caso de imágenes, el enfoque más
usado es el \emph{aprendizaje de una clase} (\textit{one-class learning}):
el modelo se entrena solo con muestras normales y asigna una puntuación de
anomalía durante la inferencia.

El detector de vector de soporte de una clase (One-Class SVM, OC-SVM) de Schölkopf et al.~\cite{scholkopf2001ocsvm} estima una
hipersuperficie que envuelve la distribución de las muestras normales en el
espacio de características. El bosque de aislamiento (Isolation
Forest) de Liu, Ting y Zhou~\cite{liu2012isolation} sigue una idea
distinta: aísla instancias construyendo árboles de decisión al azar, bajo
el supuesto de que las anomalías quedan aisladas con menos particiones que
las muestras normales. Ambos modelos exponen la interfaz
fit/predict de scikit-learn~\cite{pedregosa2011sklearn},
lo que los hace candidatos naturales para usarse con AutoPipeAD.

Métodos más recientes, como PaDiM~\cite{defard2021padim} y
PatchCore~\cite{roth2022patchcore}, aprovechan representaciones profundas de
redes neuronales preentrenadas para detectar anomalías en imágenes
industriales con alta precisión, aunque su integración es más compleja.

El dataset MVTec AD~\cite{bergmann2019mvtec} es el conjunto de referencia
más usado para evaluar métodos de detección de anomalías en imágenes
industriales. Contiene 15~categorías de objetos y texturas con defectos
reales, y se usó como referencia de compatibilidad en este trabajo.

\subsection{Preprocesamiento de Imágenes}

Gonzalez y Woods~\cite{gonzalez2018dip} describieron las transformaciones
fundamentales de imágenes digitales. Para la detección de anomalías, las
más relevantes son:

\begin{itemize}
\item \textbf{Redimensionamiento}: uniformiza las dimensiones de las
  imágenes para garantizar vectores de características de longitud fija.
\item \textbf{Suavizado gaussiano}: reduce el ruido de alta frecuencia
  convolucionando la imagen con un núcleo gaussiano de desviación
  estándar $\sigma$.
\item \textbf{Filtro de mediana}: elimina el ruido impulsivo (sal
  y pimienta) sin borrar los bordes de la imagen.
\item \textbf{Normalización estadística}: centra y escala los valores de
  intensidad según la media $\mu$ y la desviación estándar $\sigma$ del
  canal, produciendo la imagen transformada $(I - \mu)/\sigma$.
\item \textbf{CLAHE} (\textit{Contrast Limited Adaptive Histogram
  Equalization}): mejora el contraste local dividiendo la imagen en
  bloques y aplicando ecualización de histograma con un límite de
  contraste (\textit{clip limit}) que evita amplificar el ruido.
\item \textbf{Corrección gamma}: ajusta la luminosidad global mediante la
  función de potencia $I_{\text{out}} = I_{\text{in}}^{\gamma}$, donde
  $\gamma < 1$ aclara la imagen y $\gamma > 1$ la oscurece.
\item \textbf{Conversión a escala de grises}: reduce la imagen RGB a un
  solo canal de luminancia, bajando la dimensionalidad del espacio de
  características.
\end{itemize}

No existen guías generales para elegir el pipeline de preprocesamiento en
detección de anomalías: la combinación óptima depende del tipo de defecto,
las condiciones de iluminación y la arquitectura del modelo base.

\subsection{Aprendizaje Automático Automatizado y Búsqueda de Pipelines}

El aprendizaje automático automatizado (AutoML) estudia cómo
seleccionar y configurar de forma automática algoritmos y pipelines de
aprendizaje automático~\cite{yang2020automl}. La búsqueda de pipelines de
preprocesamiento es un subproblema de AutoML, donde el espacio de búsqueda
queda definido por las transformaciones disponibles y sus hiperparámetros,
discretos o continuos.

La \textbf{búsqueda exhaustiva por cuadrícula} (\textit{grid search})
evalúa el producto cartesiano de todas las opciones discretas. Garantiza
explorar el espacio completo, pero su complejidad crece como
$O\!\left(\prod_{i=1}^{n} |S_i|\right)$, donde $|S_i|$ es el número de
opciones del paso $i$ del pipeline~\cite{hastie2009elements}.

La \textbf{optimización bayesiana} construye un modelo probabilístico de la
función objetivo y lo usa para decidir los próximos puntos a evaluar,
convergiendo con menos ensayos que la búsqueda aleatoria~\cite{snoek2012practical}.
Optuna~\cite{akiba2019optuna} implementa el muestreo por Estimación de
Parzen en Árbol (Tree-structured Parzen Estimator, TPE) junto con
el podador por mediana (MedianPruner), que descarta ensayos cuyo
desempeño cae por debajo de la mediana del estudio tras un número mínimo de
evaluaciones.

\subsection{Evaluación Paralela}

Como la evaluación de cada pipeline candidato es independiente de las
demás, esta tarea se presta de forma natural al paralelismo. La librería
joblib, distribuida junto con
scikit-learn~\cite{pedregosa2011sklearn}, ofrece una interfaz de alto nivel
para paralelismo basado en procesos mediante el backend loky,
compatible con Windows, Linux y macOS sin necesidad de proteger
explícitamente el punto de entrada en scripts Python. La eficiencia de la
paralelización se mide con el speedup, definido como el cociente
$T_{\text{seq}} / T_{\text{par}}$ entre los tiempos de ejecución secuencial
y paralela.

\subsection{Métricas de Evaluación}

La evaluación de detectores de anomalías se basa en métricas de
clasificación binaria. Fawcett~\cite{fawcett2006roc} describió el área bajo
la curva ROC (AUC-ROC) como la métrica más robusta para comparar
detectores, porque no depende del umbral de decisión ni del desbalance de
clases. El AUC-ROC toma valores entre 0 y 1, donde 1.0 indica clasificación
perfecta y 0.5 corresponde a predicción aleatoria. El F1-score, la
precisión y el recall complementan el AUC-ROC al considerar un umbral de
binarización específico.

El dataset CIFAR-10~\cite{krizhevsky2009cifar}, con 60\,000 imágenes de
$32\times32$ píxeles en diez categorías, se usó como base para construir
escenarios de anomalía de una clase contra el resto
(one-vs-rest), un protocolo estándar para evaluar detectores no
supervisados sobre datos de clasificación.

% ============================================================
%  III. DESCRIPCIÓN DE LA SOLUCIÓN DESARROLLADA
% ============================================================
\section{Descripción de la Solución Desarrollada}

El repositorio de AutoPipeAD donde se encuentra la implementación completa, la documentación y los ejemplos de uso está disponible en GitHub en la dirección \url{https://github.com/ItsRichard7/AutoPipeAD}.

\subsection{Arquitectura General}

AutoPipeAD se diseñó como una librería Python con seis módulos y interfaces
públicas bien definidas. La Figura~\ref{fig:arquitectura} muestra cómo
fluyen los datos entre los componentes del sistema.

\begin{figure*}[t]
\centering
\begin{tikzpicture}[node distance=0.60cm and 0.80cm]

    %% --- Columna izquierda: actores externos ---
    \node[actor]  (usr)    {Usuario};

    %% --- Columna central: flujo principal ---
    \node[modulo, right=0.85cm of usr]  (dl)  {data\_loader\\ImageDataLoader};
    \node[modulo, below=0.60cm of dl]   (pl)  {pipeline\\BaseTransform\\Pipeline};
    \node[modulo, below=0.60cm of pl]   (se)  {search\_engine\\GridSearch\\Optuna};
    \node[modulo, below=0.60cm of se]   (ev)  {evaluator\\DataSplitter\\MetricCalc.};
    \node[modulo, below=0.60cm of ev]   (rep) {reporter\\ReportGenerator};

    %% --- Rama opcional: parallel_runner ---
    \node[modulo, right=1.5cm of se, fill=lightgray!60, draw=gray]
          (pr)  {parallel\_runner\\ParallelRunner\\(opcional)};

    %% --- Modelo base (actor externo) ---
    \node[externo, right=1.5cm of pr] (mb) {Modelo base\\(sklearn /\\PyTorch)};

    %% --- Flechas del flujo principal ---
    \draw[flecha] (usr)  -- node[above, font=\scriptsize]{dataset} (dl);
    \draw[flecha] (dl)   -- node[right, font=\scriptsize]{imágenes, etiquetas} (pl);
    \draw[flecha] (pl)   -- node[right, font=\scriptsize]{Pipeline} (se);
    \draw[flecha] (se)   -- node[right, font=\scriptsize]{usa} (ev);
    \draw[flecha] (ev)   -- node[right, font=\scriptsize]{resultados} (rep);
    \draw[flecha] (rep.west) -| (usr.south);

    %% --- Rama opcional ---
    \draw[flechaopc] (se.east) -- node[above, font=\scriptsize]{candidatos} (pr.west);
    \draw[flechaopc] (pr.south) |- node[right, font=\scriptsize]{métricas} (rep.east);

    %% --- Modelo base → módulos ---
    \draw[flechadash] (mb) -- node[above, font=\scriptsize]{fit/predict} (pr);
    \draw[flechadash] (mb.south) |- (ev.east);

\end{tikzpicture}
\caption{Arquitectura de AutoPipeAD. Las flechas con trazo continuo corresponden
  al flujo principal de búsqueda automática; las punteadas indican el módulo
  \texttt{parallel\_runner}, componente independiente y opcional.}
\label{fig:arquitectura}
\end{figure*}

\subsection{Módulo de Carga de Datos (\texttt{data\_loader})}

La clase ImageDataLoader cargó imágenes desde un directorio raíz
que debió contener los subdirectorios \texttt{normal/} (etiqueta 0) y
\texttt{anomaly/} (etiqueta 1). Se soportaron los formatos \texttt{.png},
\texttt{.jpg} y \texttt{.bmp}; los archivos con extensión no soportada se
omitieron con una advertencia de registro. El método
load() retornó un arreglo NumPy denso de forma $(n, H, W, C)$ si
todas las imágenes compartían dimensiones, o un arreglo de objetos en caso
contrario.

\subsection{Módulo de Transformaciones (\texttt{pipeline})}

\subsubsection{Clase base \texttt{BaseTransform}}

Todas las transformaciones heredan de la clase abstracta
\texttt{BaseTransform}, que define dos métodos obligatorios:

\begin{lstlisting}[language=Python]
class BaseTransform(ABC):
    @abstractmethod
    def apply(self, image: np.ndarray) -> np.ndarray:
        ...
    @abstractmethod
    def get_params(self) -> dict:
        ...
\end{lstlisting}

Este contrato permite que el motor de búsqueda trate cualquier
transformación de forma uniforme. La Figura~\ref{fig:clases} muestra la
jerarquía de clases del módulo de transformaciones.

\begin{figure*}[t]
\centering
\begin{tikzpicture}

    %% Clase base centrada
    \node[iface, minimum width=3.0cm, text width=2.8cm] at (0,0) (bt)
          {\textit{\guillemotleft{}abstracta\guillemotright{}}\\[2pt]
           \textbf{BaseTransform}\\[2pt]
           +apply(img)\\+get\_params()};

    %% Siete clases concretas, equidistantes a 2.6cm entre centros
    \node[modulo, text width=1.95cm, minimum width=2.1cm] at (-7.8,-2.6) (gb)
          {Gaussian\\Blur};
    \node[modulo, text width=1.95cm, minimum width=2.1cm] at (-5.2,-2.6) (mf)
          {Median\\Filter};
    \node[modulo, text width=1.95cm, minimum width=2.1cm] at (-2.6,-2.6) (rs)
          {Resize};
    \node[modulo, text width=1.95cm, minimum width=2.1cm] at ( 0.0,-2.6) (nm)
          {Normalize};
    \node[modulo, text width=1.95cm, minimum width=2.1cm] at ( 2.6,-2.6) (cl)
          {CLAHE\\Transform};
    \node[modulo, text width=1.95cm, minimum width=2.1cm] at ( 5.2,-2.6) (gc)
          {Gamma\\Correction};
    \node[modulo, text width=1.95cm, minimum width=2.1cm] at ( 7.8,-2.6) (gs)
          {Grayscale\\Conversion};

    %% Barra colectora horizontal (estilo bus UML)
    \draw[thick, color=tecblue] (-7.8,-1.45) -- (7.8,-1.45);

    %% Flecha de herencia: barra → BaseTransform (apunta hacia el padre)
    \draw[flecha] (0,-1.45) -- (bt.south);

    %% Stubs verticales: tope de cada nodo → barra
    \draw[thick, color=tecblue] (gb.north) -- (-7.8,-1.45);
    \draw[thick, color=tecblue] (mf.north) -- (-5.2,-1.45);
    \draw[thick, color=tecblue] (rs.north) -- (-2.6,-1.45);
    \draw[thick, color=tecblue] (nm.north) -- ( 0.0,-1.45);
    \draw[thick, color=tecblue] (cl.north) -- ( 2.6,-1.45);
    \draw[thick, color=tecblue] (gc.north) -- ( 5.2,-1.45);
    \draw[thick, color=tecblue] (gs.north) -- ( 7.8,-1.45);

\end{tikzpicture}
\caption{Jerarquía de clases del módulo \texttt{pipeline}. Las siete
  transformaciones concretas heredaron directamente de \texttt{BaseTransform}
  e implementaron \texttt{apply()} y \texttt{get\_params()}.
  La barra colectora representa herencia múltiple desde un único padre.}
\label{fig:clases}
\end{figure*}

\subsubsection{Transformaciones predefinidas}

La librería implementa siete transformaciones: GaussianBlur
(suavizado gaussiano con núcleo cuadrado de tamaño impar positivo),
MedianFilter (filtro de mediana, $\text{kernel} \geq 3$),
Normalize (normalización por media y desviación estándar, produce
float32), Resize (redimensionamiento bilineal con
OpenCV), CLAHETransform (aplica CLAHE al canal $L$ del espacio
LAB en imágenes a color, o directamente en imágenes en escala de grises),
GammaCorrection (implementada con una tabla de búsqueda para imágenes enteras) y
GrayscaleConversion (conversión de RGB a escala de grises).

\subsubsection{Clase \texttt{Pipeline}}

La clase \texttt{Pipeline} compone una secuencia ordenada de instancias de
\texttt{BaseTransform}. Aplica cada
transformación en orden, y retorna un diccionario
con clave \{índice\}\_\{Clase\} por paso, lo que permite
serializar la configuración ganadora.

\subsection{Módulo de Búsqueda (\texttt{search\_engine})}

El espacio de búsqueda se expresa como un diccionario Python donde cada
clave identifica un paso del pipeline y el valor es una lista de
transformaciones candidatas. Un ejemplo de uso real con
\texttt{val\_size=0.3}:

\begin{lstlisting}[language=Python]
search_space = {
    "resize":    [Resize(width=32, height=32)],
    "blur":      [GaussianBlur(kernel_size=3),
                  GaussianBlur(kernel_size=5)],
    "normalize": [Normalize(mean=0.5, std=0.5)],
}
engine = GridSearchEngine(
    search_space, val_size=0.3, random_state=42)
best, metrics, results = engine.run(
    images, labels,
    base_model=IsolationForest(
        contamination=0.2, random_state=0),
    metric="auc_roc")
\end{lstlisting}

Para cada paso del espacio, ambos motores añaden de forma automática y
transparente la opción de omitirlo: el usuario no configura nada
adicional y la búsqueda también evalúa pipelines sin esa transformación. Un
paso omitido simplemente no aparece en el pipeline resultante ni en su
configuración serializada, lo que preserva el contrato de tipos de
\texttt{BaseTransform} y hace trivial la reconstrucción.

\subsubsection{GridSearchEngine}

Genera el producto cartesiano de las listas de candidatos, ampliado con la opción de omisión por paso, y
evalúa cada combinación de forma exhaustiva; el espacio efectivo es de
$\prod_{i=1}^{n} (|S_i|+1)$ combinaciones. Para cada pipeline candidato:
(1)~se aplica la secuencia de transformaciones a cada imagen y se aplana el
resultado en un vector de características; (2)~se entrena una copia profunda del modelo base sobre el subconjunto de
entrenamiento; (3)~se calculan las métricas sobre el subconjunto de
validación estratificado. Se compara el valor
objetivo según la dirección de la métrica: false\_positive\_rate se
minimiza, el resto se maximiza. Los candidatos generados automáticamente que
no pueden evaluarse (por ejemplo, omitir el redimensionamiento con imágenes
de dimensiones heterogéneas) se registran con métricas NaN y una advertencia, sin
abortar la búsqueda; un candidato NaN nunca resulta ganador.

\subsubsection{OptunaSearchEngine}

Implementa optimización bayesiana con TPESampler (semilla fija) y
MedianPruner~\cite{akiba2019optuna}. Acepta el mismo formato de
espacio de búsqueda que GridSearchEngine y expone la misma firma; la omisión de cada paso se modela como un índice
categórico adicional en el muestreador. Los valores NaN que producen los
ensayos con una sola clase en el conjunto de validación (o los candidatos
inevaluables) se reemplazan por $\pm\infty$ para no corromper el estudio de
Optuna.

\subsection{Módulo de Evaluación (\texttt{evaluator})}

\subsubsection{DataSplitter}

Genera particiones train/val estratificadas de scikit-learn~\cite{pedregosa2011sklearn}. La
estratificación garantiza que la proporción de muestras anómalas sea
equivalente en ambos subconjuntos, algo importante en datasets
desbalanceados. También soporta validación cruzada $k$-pliegues. La configuración del particionador se puede guardar y
recargar en JSON o YAML.

\subsubsection{MetricCalculator}

Calcula AUC-ROC, F1-score, precisión, recall y tasa de falsos positivos.
Para obtener las puntuaciones continuas que requiere el AUC-ROC, se
prueban en orden los métodos predict\_proba,
score\_samples (negado) y decision\_function (negado)
del modelo base. Las predicciones de IsolationForest (convención
$-1$=anomalía, $+1$=normal) se normalizan a etiquetas $\{0, 1\}$.

\subsection{Módulo de Ejecución Paralela (\texttt{parallel\_runner})}

La clase \texttt{ParallelRunner} evalúa una lista de pipelines de forma
concurrente. Cada
proceso evalúa un pipeline sobre una copia independiente del modelo base,
sin estado compartido mutable. El método benchmark() mide el speedup como $T_{\text{seq}} / T_{\text{par}}$, y verifica que los resultados paralelos y
secuenciales sean numéricamente equivalentes.

\subsection{Módulo de Reportes (\texttt{reporter})}

La clase \texttt{ReportGenerator} siempre produce un reporte HTML, con exportación opcional a PDF. Si el binario no está disponible, el
sistema registra una advertencia y conserva el HTML (degradación
elegante). Las gráficas se generan con matplotlib y se incrustan como URI de datos en base64, lo que produce
un documento HTML autocontenido. La configuración del pipeline ganador se
guarda en JSON con la estructura:
\begin{lstlisting}
{
  "0_Resize":       {"width": 32, "height": 32},
  "1_GaussianBlur": {"kernel_size": 3},
  "2_Normalize":    {"mean": 0.5, "std": 0.5}
}
\end{lstlisting}
El método estático rebuild\_pipeline(config) reconstruye el
pipeline a partir de este archivo, lo que permite repetir exactamente la
misma ejecución; los pasos omitidos por la búsqueda no aparecen en la
configuración y los índices de los presentes quedan contiguos.

Si el usuario no configura un directorio de salida, los archivos se
escriben en una subcarpeta reports/ 
junto al script que invoca la librería; en sesiones interactivas (Jupyter)
se usa el directorio de trabajo, que coincide con el del cuaderno. La
carpeta se crea automáticamente, de modo que el caso de uso estándar no
requiere configurar rutas.

\subsection{Flujo de Evaluación de un Pipeline}

La Figura~\ref{fig:flujo} muestra el diagrama de actividad del proceso de
evaluación de un pipeline candidato dentro del motor de búsqueda.

\begin{figure}[t]
\centering
\begin{tikzpicture}[node distance=0.75cm]
    \node[startend] (ini)    {};
    \node[estado,   below=0.55cm of ini] (load)   {Cargar imagen};
    \node[decision, below=of load]        (valid)  {\scriptsize¿Formato soportado?};
    \node[estado,   right=1.6cm of valid] (omit)   {Omitir +\\advertencia};
    \node[estado,   below=of valid]       (trans)  {Aplicar transformación $i$};
    \node[decision, below=of trans]       (more)   {\scriptsize¿Más pasos?};
    \node[estado,   below=of more]        (score)  {Puntuar con\\modelo base};
    \node[decision, below=of score]       (imgs)   {\scriptsize¿Más imágenes?};
    \node[estado,   below=of imgs]        (metric) {Calcular métricas\\(AUC-ROC, F1\ldots)};
    \node[startend, below=0.55cm of metric] (fin)  {};

    \draw[flecha] (ini)    -- (load);
    \draw[flecha] (load)   -- (valid);
    \draw[flecha] (valid)  -- node[right, font=\scriptsize]{sí} (trans);
    \draw[flecha] (valid)  -- node[above, font=\scriptsize]{no} (omit);
    \draw[flecha] (trans)  -- (more);
    \draw[flecha] (more.west) to[out=180,in=180]
                 node[left, font=\scriptsize]{sí,\,$i{+}1$} (trans.west);
    \draw[flecha] (more)   -- node[right, font=\scriptsize]{no} (score);
    \draw[flecha] (score)  -- (imgs);
    \draw[flecha] (imgs.east) to[out=0,in=0]
                 node[right, font=\scriptsize]{sí} (load.east);
    \draw[flecha] (imgs)   -- node[right, font=\scriptsize]{no} (metric);
    \draw[flecha] (metric) -- (fin);
\end{tikzpicture}
\caption{Diagrama de actividad de la evaluación de un pipeline candidato.
  Este flujo se ejecutó para cada combinación generada por el motor de búsqueda.}
\label{fig:flujo}
\end{figure}

\subsection{Decisiones de Diseño}

\textbf{Inyección del modelo base}: el motor de búsqueda recibe el modelo
base como argumento y lo clona antes de cada
evaluación. Así se garantiza que todos los candidatos se evalúen sobre un
modelo sin entrenar idéntico, preservando la comparabilidad y la
reproducibilidad.

\textbf{Funciones a nivel de módulo para lógica compartida}: las funciones
\texttt{evaluate\_pipeline()}, \texttt{is\_better()} y
\texttt{validate\_metric()} se definieron a nivel de módulo, de modo que
\texttt{GridSearchEngine}, \texttt{OptunaSearchEngine} y
\texttt{ParallelRunner} pudieran usarlas sin necesidad de herencia.

\textbf{Elección de joblib sobre multiprocessing}: se eligió joblib porque abstrae las diferencias entre plataformas
(modelo fork en Linux versus spawn en Windows), es
compatible con ambos motores de búsqueda y es la misma dependencia que usa
scikit-learn~\cite{pedregosa2011sklearn}.

\textbf{Elección de Jinja2 + pdfkit}: se descartó WeasyPrint porque
requiere la pila gráfica GTK (Pango/Cairo), no disponible en Windows sin
configuración compleja. En su lugar se adoptó pdfkit con
tolerancia a fallos.

\textbf{Omisión de pasos por ausencia}: se eligió representar el paso
omitido como ausencia (no aparece en el pipeline ni en su configuración
serializada) en lugar de un marcador explícito o una transformación
identidad: así se preserva el contrato de tipos de \texttt{BaseTransform}
y la reconstrucción desde JSON/YAML funciona sin casos especiales,
incluido el pipeline vacío.

\subsection{Tecnologías Utilizadas}

La Tabla~\ref{tab:tecnologias} resume las dependencias principales.

\begin{table}[t]
\centering
\renewcommand{\arraystretch}{1.2}
\caption{Dependencias principales de AutoPipeAD.}
\label{tab:tecnologias}
\begin{tabular}{@{}lll@{}}
\toprule
\textbf{Componente} & \textbf{Librería} & \textbf{Versión mínima} \\
\midrule
Procesamiento de imágenes & OpenCV        & 4.9  \\
                          & scikit-image  & 0.22 \\
                          & Pillow        & 10.0 \\
Cómputo numérico          & NumPy         & 1.26 \\
Modelos ML e interfaz     & scikit-learn  & 1.4  \\
                          & PyTorch       & 2.2  \\
Búsqueda bayesiana        & Optuna        & 3.6  \\
Paralelismo               & joblib        & 1.3  \\
Visualización             & matplotlib    & 3.8  \\
Plantillas HTML           & Jinja2        & 3.1  \\
Exportación PDF           & pdfkit        & 1.0  \\
Serialización             & PyYAML        & 6.0  \\
\bottomrule
\end{tabular}
\end{table}

El entorno de desarrollo utilizó Python~3.12 con las herramientas Ruff
(verificador de estilo e importaciones) y mypy (verificación estática de
tipos estricta).

% ============================================================
%  IV. RESULTADOS
% ============================================================
\section{Resultados}

\subsection{Suite de Pruebas}

Se desarrolló una suite de pruebas con \textit{pytest}, compuesta por seis
archivos de pruebas unitarias e integración —uno por módulo— y un archivo
de pruebas de extremo a extremo que
verifica el flujo completo, desde la carga del dataset hasta la generación
del reporte y la reconstrucción del pipeline ganador: 134~pruebas en total
con una cobertura del 97\,\% del código fuente según pytest-cov.

\subsection{Experimento 1: Dataset Sintético con Grid Search y Optuna}

El primer notebook ejecutó la búsqueda sobre un
dataset sintético de 50~imágenes de $64\times64\times3$ píxeles, generado
de forma programática: las 40~imágenes normales se crearon con valores en
el rango $[0, 40]$ (oscuras) y las 10~anómalas con valores en $[200, 256]$
(brillantes), de modo que las anomalías representan un 20\,\% del conjunto,
acorde con el supuesto de rareza de los detectores de anomalías. El modelo
base fue \texttt{IsolationForest} de
scikit-learn~\cite{pedregosa2011sklearn} con \texttt{contamination=0.2}
(igual a la proporción real de anomalías) y \texttt{random\_state=0}. El
espacio de búsqueda y los parámetros de evaluación se muestran en la
Tabla~\ref{tab:espacio_exp1}.

\begin{table}[t]
\centering
\renewcommand{\arraystretch}{1.2}
\caption{Configuración del Experimento~1 (dataset sintético).}
\label{tab:espacio_exp1}
\begin{tabular}{@{}ll@{}}
\toprule
\textbf{Parámetro} & \textbf{Valor} \\
\midrule
Imágenes totales    & 50 (40 normales, 10 anómalas) \\
Dimensiones         & $64\times64\times3$ \\
\texttt{val\_size}  & 0.30 \\
\texttt{random\_state} & 42 \\
Modelo base         & IsolationForest (\texttt{contam=0.2}, \texttt{random\_st=0}) \\
Métrica objetivo    & AUC-ROC \\
Paso \emph{resize}  & 1 opción: Resize($32\times32$) \\
Paso \emph{blur}    & 2 opciones: GaussianBlur($k=3$), GaussianBlur($k=5$) \\
Paso \emph{normalize} & 1 opción: Normalize($\mu=0.5$, $\sigma=0.5$) \\
Total (Grid Search) & 12 combinaciones ($2\cdot3\cdot2$, con omisión por paso) \\
Ensayos (Optuna)    & 8 \\
\bottomrule
\end{tabular}
\end{table}

La búsqueda exhaustiva evaluó las 12~combinaciones y todas —incluidos los
subconjuntos generados por la omisión automática y el pipeline vacío—
alcanzaron métricas perfectas (Tabla~\ref{tab:metricas_exp1}): la
separación de brillo entre clases es tan marcada que cualquier
preprocesamiento (o ninguno) preserva la separabilidad, y el umbral
calibrado con \texttt{contamination=0.2} produce predicciones binarias
correctas. El ganador fue el primer candidato en alcanzar el óptimo según
el orden de exploración: Resize($32\times32$) $\to$
GaussianBlur($k=3$) $\to$ Normalize.

\begin{table}[t]
\centering
\renewcommand{\arraystretch}{1.2}
\caption{Resultados de Grid Search — Experimento~1 (extracto; las
  12~combinaciones obtuvieron métricas idénticas).}
\label{tab:metricas_exp1}
\begin{tabular}{@{}lccccc@{}}
\toprule
\textbf{Pipeline} & \textbf{AUC-ROC} & \textbf{F1} & \textbf{Prec.} &
\textbf{Recall} & \textbf{FPR} \\
\midrule
\#0 Res+Blur+Norm (\textbf{winner}) & 1.0 & 1.0 & 1.0 & 1.0 & 0.0 \\
\#11 (pipeline vacío)                            & 1.0 & 1.0 & 1.0 & 1.0 & 0.0 \\
\bottomrule
\end{tabular}
\end{table}

Durante el desarrollo se observó que una versión balanceada 50/50 de este
mismo dataset degradaba el AUC-ROC del mejor pipeline a 0.5000
(predicción aleatoria): IsolationForest asume que las anomalías son raras
y, con la clase anómala representando la mitad del entrenamiento, el
modelo llegó a clasificar todas las muestras de validación como anómalas
(FPR=1.0). Por ello el experimento final usa una proporción realista de
20\,\% de anomalías y el parámetro \texttt{contamination=0.2}, que calibra
el umbral de decisión binaria a la tasa de anomalías esperada.

La búsqueda bayesiana con Optuna ejecutó 8~ensayos y alcanzó el mismo
AUC-ROC de 1.0000 con un pipeline distinto pero igualmente óptimo
(GaussianBlur($k=3$) $\to$ Normalize, omitiendo el redimensionamiento):
cuando varios candidatos comparten el valor óptimo, cada motor retiene el
primero que lo alcanza según su propio orden de exploración.

La reproducibilidad se verificó así: la configuración del pipeline
ganador se guardó automáticamente, y el método rebuild\_pipeline() reconstruyó exactamente el mismo pipeline,
confirmado al comparar los parámetros de ambos.

\subsection{Experimento 2: CIFAR-10 Sintético con Optuna}

El segundo notebook construyó un
escenario de anomalía de una clase contra el resto sobre un dataset
CIFAR-10 sintético~\cite{krizhevsky2009cifar}: las imágenes normales
(clase 1) se generaron con tonos verdosos ($\mu=[60,150,60]$,
$\sigma=12$) y las anómalas (demás clases) con tonos rojizos
($\mu=[170,60,60]$, $\sigma=40$). La Tabla~\ref{tab:espacio_exp2} muestra
la configuración del experimento.

\begin{table}[t]
\centering
\renewcommand{\arraystretch}{1.2}
\caption{Configuración del Experimento~2 (CIFAR-10 sintético).}
\label{tab:espacio_exp2}
\begin{tabular}{@{}ll@{}}
\toprule
\textbf{Parámetro} & \textbf{Valor} \\
\midrule
Imágenes totales    & 100 (60 normales, 40 anómalas) \\
Dimensiones         & $32\times32\times3$ \\
\texttt{val\_size}  & 0.30 \\
\texttt{random\_state} & 42 \\
Motor               & Optuna (12 ensayos, TPE + MedianPruner) \\
Paso \emph{resize}  & 1 opción: Resize($16\times16$) \\
Paso \emph{median}  & 1 opción: MedianFilter($k=3$) \\
Paso \emph{blur}    & 2 opciones: GaussianBlur($k=3$), GaussianBlur($k=5$) \\
Paso \emph{clahe}   & 1 opción: CLAHE (clip $2.0$, rejilla $4\times4$) \\
Paso \emph{gamma}   & 2 opciones: $\gamma=0.5$, $\gamma=1.5$ \\
Paso \emph{gray}    & 1 opción: GrayscaleConversion \\
Paso \emph{normalize} & 1 opción: Normalize($\mu=0.5$, $\sigma=0.5$) \\
Espacio total       & 288 combinaciones (con omisión por paso) \\
\bottomrule
\end{tabular}
\end{table}

Este experimento usa un espacio de búsqueda completo con las siete
transformaciones predefinidas de la librería. De los 12~ensayos de Optuna,
11 alcanzaron AUC-ROC=1.0000 (el restante, 0.9954). El pipeline ganador
correspondió al primer ensayo evaluado: MedianFilter($k=3$) $\to$
GaussianBlur($k=3$) $\to$ CLAHE($2.0$) $\to$ GrayscaleConversion
(Tabla~\ref{tab:metricas_exp2}), un pipeline que \emph{omite} tres de los
siete pasos disponibles (redimensionamiento, corrección gamma y
normalización), lo que ilustra la exploración automática de subconjuntos
del espacio de búsqueda.

\begin{table}[t]
\centering
\renewcommand{\arraystretch}{1.2}
\caption{Métricas del pipeline ganador — Experimento~2 (CIFAR-10 sintético).}
\label{tab:metricas_exp2}
\begin{tabular}{@{}lc@{}}
\toprule
\textbf{Métrica} & \textbf{Valor} \\
\midrule
AUC-ROC              & 1.0 \\
F1-score             & 1.0 \\
Precisión            & 1.0 \\
Recall               & 1.0 \\
Tasa falsos positivos& 0.0 \\
\bottomrule
\end{tabular}
\end{table}

El AUC-ROC=1.0000 refleja la alta separabilidad del dataset sintético: la
diferencia cromática entre la clase normal (verde) y la anómala (rojo) fue
suficiente para que IsolationForest separara perfectamente ambas clases
con casi cualquier subconjunto de transformaciones. Este resultado
confirma la correctitud funcional del sistema sobre un espacio amplio
(288 combinaciones exploradas con solo 12 ensayos) y muestra que la
omisión automática de pasos permite encontrar pipelines más cortos con el
mismo desempeño.

\subsection{Experimento 3: Adaptador para MVTec AD}

El tercer notebook  mostró la integración
de AutoPipeAD con el dataset MVTec AD~\cite{bergmann2019mvtec}, cuya
estructura de directorios (categorías con subdirectorios por tipo de
defecto) es distinta del formato que espera \texttt{ImageDataLoader}. El
notebook incluye la función \texttt{build\_normal\_anomaly\_dir()}, que
convierte la estructura \texttt{test/good/} + \texttt{test/\textit{defecto}/}
al formato \texttt{normal/} + \texttt{anomaly/}. Como no se contó con el
dataset real (requiere registro y descarga de
$\sim$4.9~GB), el notebook se probó con un dataset sintético de
validación (30~imágenes normales con textura gris uniforme y
18~anómalas con parche brillante), lo que confirmó que el adaptador
funciona correctamente. El espacio de búsqueda usó \texttt{CLAHETransform}
con $\text{clip\_limit} \in \{2.0, 4.0\}$ y \texttt{GaussianBlur} con
$k \in \{3, 5\}$, junto con \texttt{Resize}($64\times64$) y
\texttt{Normalize}; con la omisión automática por paso, la búsqueda
exhaustiva evaluó 36~combinaciones ($2\cdot3\cdot3\cdot2$). De ellas,
24 alcanzaron AUC-ROC=1.0000 y la peor obtuvo 0.8704; el ganador fue
Resize($64\times64$) $\to$ CLAHE($2.0$) $\to$ GaussianBlur($k=3$) $\to$
Normalize, con las cinco métricas perfectas.

\subsection{Evaluación del Módulo de Paralelismo}

El método \texttt{ParallelRunner.benchmark()} ejecutó la evaluación del
espacio de búsqueda del Experimento~1 (12~pipelines, incluida la omisión
por paso) en modo secuencial (\texttt{n\_jobs=1}) y en modo paralelo
(\texttt{n\_jobs=-1}), obteniendo los resultados de la
Tabla~\ref{tab:speedup}.

\begin{table}[t]
\centering
\renewcommand{\arraystretch}{1.2}
\caption{Resultados de la evaluación de paralelismo — Experimento~1.}
\label{tab:speedup}
\begin{tabular}{@{}lcc@{}}
\toprule
\textbf{Parámetro} & \textbf{Valor medido} \\
\midrule
\textit{Speedup} medido              & 0.25$\times$ \\
Consistencia (seq. vs. par.)         & Verdadero \\
\textit{n\_jobs}                     & $-1$ (todos los núcleos) \\
Pipelines evaluados                  & 12 \\
\bottomrule
\end{tabular}
\end{table}

El \textit{speedup} de 0.25$\times$ (la ejecución paralela fue alrededor
de 4~veces más lenta que la secuencial) es el comportamiento esperado
para un espacio de búsqueda pequeño de evaluación barata: la sobrecarga
de inicializar el \textit{pool} de procesos \textit{loky} domina sobre el
tiempo de cómputo de cada pipeline, aunque la penalización se reduce al
crecer el espacio (con 2~pipelines, en mediciones previas, el
\textit{speedup} fue de 0.06$\times$). El beneficio de la paralelización
aparece en espacios más amplios o con modelos base costosos, donde el
tiempo de evaluación supera esa sobrecarga. El indicador
\texttt{consistent=True} confirma que los resultados paralelos y
secuenciales fueron numéricamente equivalentes, lo que valida la
correctitud del mecanismo.

\subsection{Comparación de Estrategias de Búsqueda}

La Tabla~\ref{tab:comparacion} compara las dos estrategias de búsqueda en
el Experimento~1, donde el espacio fue lo bastante pequeño para evaluar
ambas sobre el mismo conjunto.

\begin{table}[t]
\centering
\renewcommand{\arraystretch}{1.2}
\caption{Comparación Grid Search vs.\ Optuna — Experimento~1.}
\label{tab:comparacion}
\begin{tabular}{@{}lccc@{}}
\toprule
\textbf{Estrategia} & \textbf{Evaluaciones} & \textbf{AUC-ROC ganador} & \textbf{Pipeline ganador} \\
\midrule
Grid Search & 12 & 1.0000 & Resize + Blur + Norm \\
Optuna      & 8  & 1.0000 & Blur + Norm \\
\bottomrule
\end{tabular}
\end{table}

Ambas estrategias alcanzaron el mismo AUC-ROC ganador, y Optuna lo hizo
con menos evaluaciones (8 frente a las 12 del espacio completo). Las
composiciones ganadoras difieren porque en este dataset varios candidatos
comparten el valor óptimo y cada motor retiene el primero que lo alcanza
según su orden de exploración; ambas son igualmente válidas. La ventaja de
la optimización bayesiana se acentúa en espacios con cientos de
combinaciones, como el del Experimento~2 (288 combinaciones exploradas con
12 ensayos).

\subsection{Fortalezas de la Solución}

\begin{enumerate}
\item \textbf{Reproducibilidad garantizada}: la semilla fija en el
  particionador y en el muestreador de Optuna
  produce resultados idénticos en ejecuciones independientes, verificado
  con el pipeline reconstruido.
\item \textbf{Exploración automática de subconjuntos}: cada paso del
  espacio incluye la opción transparente de omitirse, sin configuración
  adicional del usuario; la búsqueda puede así descartar transformaciones
  que no aportan (el ganador del Experimento~2 omitió tres de los siete
  pasos disponibles).
\item \textbf{Extensibilidad}: la interfaz \texttt{BaseTransform} permite
  agregar transformaciones personalizadas sin modificar ningún otro
  módulo.
\item \textbf{Compatibilidad con la interfaz estándar}: la librería acepta
  cualquier modelo con métodos fit/predict de
  scikit-learn, así como modelos PyTorch con la misma interfaz.
\item \textbf{API concisa}: el flujo completo de carga, búsqueda y
  reporte se expresa en menos de 20~líneas de código de usuario, según
  muestran los notebooks de ejemplo.
\item \textbf{Reporte autocontenido y de ubicación predecible}: el
  documento HTML incluye gráficas e imágenes de anomalías codificadas en
  base64, sin depender de archivos externos; sin configuración, los
  archivos se escriben junto al script o cuaderno
  invocador, con la carpeta creada automáticamente.
\item \textbf{Adaptador de estructura de datos}: Se demuestra que la librería puede
  integrarse con datasets de formato distinto mediante adaptadores
  externos.
\end{enumerate}

\subsection{Limitaciones de la Solución}

\begin{enumerate}
\item \textbf{Sensibilidad al balance de clases}: IsolationForest asume
  que las anomalías son raras; en pruebas de desarrollo con un balance
  50/50 en el entrenamiento, el AUC-ROC bajó a 0.5000 (predicción
  aleatoria). El Experimento~1 final mitiga esta limitación del modelo
  base con una proporción realista (20\,\% de anomalías) y el parámetro contamination acorde.
\item \textbf{Paralelismo no beneficioso para espacios pequeños}: el
  \textit{speedup} de 0.25$\times$ con 12~pipelines de evaluación barata
  muestra que la paralelización resulta contraproducente cuando la
  sobrecarga de crear los procesos supera el costo de cada evaluación.
\item \textbf{Representación por aplanamiento de píxeles}: la extracción
  de características aplana la imagen en un vector, lo que descarta la
  estructura espacial local. Esto limita la compatibilidad con modelos
  basados en características profundas, como PatchCore o PaDiM.
\item \textbf{Validación simple}: la búsqueda usa una sola partición
  train/val, lo que introduce varianza en la estimación
  de la métrica para datasets pequeños. La validación cruzada $k$-pliegues
  ya está implementada, pero todavía no se integró
  en el bucle de búsqueda automática.
\item \textbf{Exportación a PDF condicional}: generar el PDF requiere el
  binario externo wkhtmltopdf, que no viene incluido entre las
  dependencias de la librería.
\end{enumerate}

% ============================================================
%  V. CONCLUSIONES
% ============================================================
\section{Conclusiones}

AutoPipeAD implementa un sistema funcional y reproducible para seleccionar
automáticamente pipelines de preprocesamiento en detección de anomalías en
imágenes. La arquitectura modular de seis componentes, con interfaces bien
definidas, separa las responsabilidades de carga de datos, definición de
transformaciones, búsqueda de configuraciones, evaluación de métricas,
ejecución paralela y generación de reportes, lo que permite extender la
librería sin modificar componentes existentes.

Los experimentos confirmaron la correctitud funcional del sistema. En el
Experimento~1, las dos estrategias de búsqueda alcanzaron un AUC-ROC de
1.0000 de forma determinista bajo una semilla fija: Grid Search evaluó las 12~combinaciones del
espacio (incluidos los subconjuntos generados por la omisión automática de
pasos) y Optuna llegó al mismo óptimo con solo 8~ensayos. Reconstruir el
pipeline desde el archivo JSON generado dio resultados idénticos, y las
métricas de los tres notebooks resultaron reproducibles entre ejecuciones
independientes. En el Experimento~2, la búsqueda bayesiana identificó
desde el primer ensayo un pipeline con AUC-ROC de 1.0000 sobre un espacio
de 288~combinaciones con las siete transformaciones predefinidas,
omitiendo tres pasos que no aportaban al desempeño.

La evaluación paralela mediante benchmark para
12~pipelines dio un de 0.25$\times$, lo que muestra que
la sobrecarga del backend superó el beneficio computacional
en espacios de búsqueda de evaluación barata. La consistencia entre
resultados paralelos y secuenciales se verificó en todos los experimentos.

El hallazgo más relevante durante el desarrollo fue la sensibilidad de
IsolationForest al balance de clases en el conjunto de entrenamiento: con
una proporción 50/50 de muestras normales y anómalas, el detector operó al
nivel de una predicción aleatoria (AUC-ROC=0.5). Esta limitación es
propia del modelo base, no del sistema de selección de pipelines, y se
mitigó en los experimentos finales usando una proporción realista de
anomalías (20\,\%) junto con el parámetro contamination
calibrado a esa tasa.

% ============================================================
%  VI. RECOMENDACIONES Y TRABAJO FUTURO
% ============================================================
\section{Recomendaciones y Trabajo Futuro}

\subsection{Integración de Validación Cruzada en el Bucle de Búsqueda}

Incorporar la validación cruzada $k$-pliegues (ya implementada en
\texttt{DataSplitter}) dentro del bucle de búsqueda automática reduciría
la varianza al estimar métricas en datasets pequeños. La métrica
objetivo podría promediarse sobre los $k$~pliegues en lugar de evaluarse
sobre una única partición.

\subsection{Extracción de Características Profundas}

Agregar un extractor de características basado en redes neuronales
convolucionales preentrenadas —como ResNet o los
descriptores de parches de PatchCore como etapa
adicional del pipeline aumentaría la capacidad discriminativa sin
necesitar etiquetas de anomalías durante el entrenamiento, superando la
limitación del aplanamiento de píxeles.

\subsection{Soporte de Hiperparámetros Continuos}

Los hiperparámetros continuos del preprocesamiento —como el valor
$\sigma$ del suavizado gaussiano o el exponente $\gamma$— podrían
incorporarse al espacio de búsqueda de Optuna como variables reales, aprovechando la capacidad nativa del
muestreador TPE para espacios mixtos.

\subsection{Interfaz de Línea de Comandos}

Desarrollar una interfaz de línea de comandos (CLI) que lea la
configuración desde el archivo YAML de ejemplo permitiría ejecutar búsquedas completas sin escribir código Python,
ampliando el acceso a la herramienta.

\subsection{Validación con Datasets Reales de Anomalías Industriales}

Los experimentos de este trabajo usaron datasets sintéticos pensados para
validar la correctitud funcional. Validar con datos reales de MVTec
AD —con defectos industriales sutiles— permitiría
medir el aporte real del preprocesamiento automático sobre modelos de
referencia como IsolationForest u OC-SVM.

\subsection{Comparación con Herramientas AutoML Existentes}

Comparar empíricamente el rendimiento y el tiempo de búsqueda de
AutoPipeAD frente a herramientas AutoML generalistas como
Auto-sklearn daría evidencia del valor diferencial
de especializarse en detección de anomalías.

% ============================================================
%  BIBLIOGRAFÍA
% ============================================================
\bibliographystyle{IEEEtran}
\begin{thebibliography}{99}

\bibitem{chandola2009anomaly}
V.~Chandola, A.~Banerjee, y V.~Kumar,
``Anomaly detection: A survey,''
\emph{ACM Computing Surveys}, vol.~41, no.~3, pp.~1--58, sep.~2009.

\bibitem{bergmann2019mvtec}
P.~Bergmann, M.~Fauser, D.~Sattlegger, y C.~Steger,
``MVTec AD -- A comprehensive real-world dataset for unsupervised anomaly
  detection,''
en \emph{Proc. IEEE Conf. Comput. Vis. Pattern Recognit. (CVPR)},
Long Beach, CA, EE.~UU., jun.~2019, pp.~9592--9600.

\bibitem{roth2022patchcore}
K.~Roth, L.~Pemula, J.~Zepeda, B.~Schölkopf, T.~Brox, y P.~Gehler,
``Towards total recall in industrial anomaly detection,''
en \emph{Proc. IEEE Conf. Comput. Vis. Pattern Recognit. (CVPR)},
Nueva Orleans, LA, EE.~UU., jun.~2022, pp.~14318--14328.

\bibitem{defard2021padim}
T.~Defard, A.~Setkov, A.~Loesch, y R.~Audigier,
``PaDiM: A patch distribution modeling framework for anomaly detection and
  localization,''
en \emph{Proc. Int. Conf. Pattern Recognit. (ICPR)},
Milán, Italia, ene.~2021, pp.~475--489.

\bibitem{liu2012isolation}
F.~T.~Liu, K.~M.~Ting, y Z.-H.~Zhou,
``Isolation-based anomaly detection,''
\emph{ACM Trans. Knowl. Discov. Data}, vol.~6, no.~1, pp.~3:1--3:39, mar.~2012.

\bibitem{scholkopf2001ocsvm}
B.~Schölkopf, J.~C.~Platt, J.~Shawe-Taylor, A.~J.~Smola,
y R.~C.~Williamson,
``Estimating the support of a high-dimensional distribution,''
\emph{Neural Computation}, vol.~13, no.~7, pp.~1443--1471, jul.~2001.

\bibitem{akiba2019optuna}
T.~Akiba, S.~Sano, T.~Yanase, T.~Ohta, y M.~Koyama,
``Optuna: A next-generation hyperparameter optimization framework,''
en \emph{Proc. 25th ACM SIGKDD Int. Conf. Knowl. Discov. Data Mining
  (KDD)}, Anchorage, AK, EE.~UU., ago.~2019, pp.~2623--2631.

\bibitem{snoek2012practical}
J.~Snoek, H.~Larochelle, y R.~P.~Adams,
``Practical Bayesian optimization of machine learning algorithms,''
en \emph{Adv. Neural Inf. Process. Syst. (NeurIPS)}, vol.~25,
Lake Tahoe, NV, EE.~UU., dic.~2012, pp.~2951--2959.

\bibitem{yang2020automl}
L.~Yang y A.~Shami,
``On hyperparameter optimization of machine learning algorithms: Theory and
  practice,''
\emph{Neurocomputing}, vol.~415, pp.~295--316, nov.~2020.

\bibitem{gonzalez2018dip}
R.~C.~Gonzalez y R.~E.~Woods,
\emph{Digital Image Processing}, 4.ª~ed.
Hoboken, NJ, EE.~UU.: Pearson, 2018.

\bibitem{hastie2009elements}
T.~Hastie, R.~Tibshirani, y J.~Friedman,
\emph{The Elements of Statistical Learning: Data Mining, Inference, and
  Prediction}, 2.ª~ed.
Nueva York, NY, EE.~UU.: Springer, 2009.

\bibitem{pedregosa2011sklearn}
F.~Pedregosa, G.~Varoquaux, A.~Gramfort, V.~Michel, B.~Thirion,
O.~Grisel, M.~Blondel, P.~Prettenhofer, R.~Weiss, V.~Dubourg,
J.~Vanderplas, A.~Passos, D.~Cournapeau, M.~Brucher, M.~Perrot,
y E.~Duchesnay,
``Scikit-learn: Machine learning in {Python},''
\emph{J. Mach. Learn. Res.}, vol.~12, pp.~2825--2830, oct.~2011.

\bibitem{fawcett2006roc}
T.~Fawcett,
``An introduction to ROC analysis,''
\emph{Pattern Recognit. Lett.}, vol.~27, no.~8, pp.~861--874, jun.~2006.

\bibitem{krizhevsky2009cifar}
A.~Krizhevsky,
``Learning multiple layers of features from tiny images,''
Univ. Toronto, Toronto, ON, Canadá, Inf.~Téc., 2009.

\bibitem{he2016resnet}
K.~He, X.~Zhang, S.~Ren, y J.~Sun,
``Deep residual learning for image recognition,''
en \emph{Proc. IEEE Conf. Comput. Vis. Pattern Recognit. (CVPR)},
Las Vegas, NV, EE.~UU., jun.~2016, pp.~770--778.

\end{thebibliography}

\end{document}